\appendix
\chapter{Reproducibility Commands}
This appendix lists command templates only. Methodology, metric definitions, and interpretation are covered in Chapter 4.
For portability and privacy, local absolute paths are replaced by placeholders.
\texttt{<RAW\_DIR\_X>} and \texttt{<FILTERED\_DIR\_X>} denote dataset-specific input roots for dataset \(X\); \texttt{<DB\_X>} denotes an output databank path.

\section{Filtered-Input Benchmark Commands (Primary Evaluation)}
Use one command per dataset:
\begin{lstlisting}[breaklines=true,basicstyle=\ttfamily\small]
python3 process_dicom.py <FILTERED_DIR_X> <DB_X> --no-subdirs --timing --max-workers 8
\end{lstlisting}

\section{Three-Stage Size Measurement Commands}
For each dataset, measure:
\begin{enumerate}
\item raw input folder size,
\item one-instance-per-series filtered folder size,
\item processed databank size.
\end{enumerate}

One-instance-per-series subsets were generated with:
\begin{lstlisting}[breaklines=true,basicstyle=\ttfamily\small]
python3 extract_one_per_series.py <RAW_DIR_X> <ONE_PER_SERIES_DIR_X>
\end{lstlisting}

Size measurement (bytes):
\begin{lstlisting}[breaklines=true,basicstyle=\ttfamily\small]
du -sb <RAW_DIR_X> <ONE_PER_SERIES_DIR_X> <DB_X>
\end{lstlisting}

\section{Raw vs Filtered Time and Memory Commands}
\begin{lstlisting}[breaklines=true,basicstyle=\ttfamily\small]
/usr/bin/time -l python3 process_dicom.py <RAW_DIR_X> <RAW_DB_X> --no-subdirs --timing --max-workers 8
/usr/bin/time -l python3 process_dicom.py <FILTERED_DIR_X> <FILTERED_DB_X> --no-subdirs --timing --max-workers 8
\end{lstlisting}

\section{Repeated Runtime Stability Commands}
Run five repetitions per dataset:
\begin{lstlisting}[breaklines=true,basicstyle=\ttfamily\small]
python3 process_dicom.py <FILTERED_DIR_X> <REPEAT_DB_X> --no-subdirs --timing --no-auto-workers --max-workers 8
\end{lstlisting}

\section{Automated Test Execution Commands}
The implemented automated test suite can be executed locally with:
\begin{lstlisting}[breaklines=true,basicstyle=\ttfamily\small]
python3 -m unittest discover -s tests -v
\end{lstlisting}

The suite currently contains 58 tests spanning discovery, extraction helpers, processing helpers, storage integration, and anonymized export/web utility behavior.

\section{Continuous Integration Configuration}
Regression checks are also executed in GitHub Actions using:
\begin{lstlisting}[breaklines=true,basicstyle=\ttfamily\small]
.github/workflows/tests.yml
\end{lstlisting}

The workflow runs on push and pull-request events, installs project requirements, and executes the same unittest discovery command.

\chapter{Supplementary Evaluation Notes}
Values in Chapter 4 are environment-specific baselines and should be interpreted comparatively. For publication-grade benchmarking, keep controlled system load, fixed worker settings, and repeated measurements with variance reporting.

\section{Hardware Specification}
All benchmarks in this thesis were executed on the local development machine used throughout the project:
\begin{itemize}
\item Model: MacBook Pro (Model Identifier: Mac15,3)
\item Chip: Apple M3
\item CPU Cores: 8 total (4 performance + 4 efficiency)
\item Memory: 8 GB
\item Architecture: arm64
\item Operating System: macOS 26.3 (Build 25D125)
\end{itemize}

\section{Project Material on GitHub}
All project materials used in this thesis are maintained in the GitHub repository:
\begin{center}
\url{https://github.com/cccelik/DICOM-Metadata-Browser}
\end{center}

The repository contains:
\begin{itemize}
\item source modules such as \path{process_dicom.py}, \path{extract_metadata.py}, \path{store_metadata.py}, and \path{webui.py},
\item the automated test suite under \path{tests/},
\item CI workflow definitions under \path{.github/workflows/}.
\item the thesis report sources and generated document under \texttt{Efficient Data Access and Storage Optimization for}\linebreak[1]\texttt{ PET:CT Imaging Data/}.
\end{itemize}
